% Part: incompleteness
% Chapter: incompleteness-provability
% Section: second-incompleteness-thm

\documentclass[../../../include/open-logic-section]{subfiles}

\begin{document}

\olfileid[fr]{inc}{inp}{2in}

\olsection{Le second théorème d'incomplétude}

Comment exprimer l'affirmation que $\Th{PA}$ ne
démontre pas sa propre cohérence? Dire que $\Th{PA}$
est incohérente revient à dire qu'elle démontre
une contradiction, par exemple
$\Th{PA}\Proves\eq[0][1]$. Nous prendrons pour
énoncé de cohérence $\OCon[\Th{PA}]$ l'!!{sentence}
$\lnot\OProv[\Th{PA}](\gn{\lfalse})$.
Ces deux témoins de contradiction donnent des énoncés de
cohérence équivalents dans~$\Th{PA}$. La théorie démontre
$\eq[0][1]\lif\lfalse$, puisque $\Th{Q}$ démontre que les
numéraux de zéro et de un sont distincts, et elle démontre
$\lfalse\lif\eq[0][1]$ par logique. En appliquant P1,
puis P2, à ces deux implications, elle obtient

\[
\OProv[\Th{PA}](\gn{\eq[0][1]})
\liff\OProv[\Th{PA}](\gn{\lfalse}).
\]
Leurs négations sont donc également équivalentes dans la
théorie. Le théorème suivant exprime le résultat recherché.

\begin{thm}
\ollabel{thm:second-incompleteness}
Si $\Th{PA}$ est cohérente, elle ne démontre pas
$\OCon[\Th{PA}]$.
\end{thm}

Il importe de noter que le théorème dépend de la
représentation particulière de
$\OCon[\Th{PA}]$, c'est-à-dire de celle de
$\OProv[\Th{PA}](y)$. Nous n'utiliserons que les
trois conditions de dérivabilité satisfaites par
cette dernière. Le résultat s'étend donc à toute
théorie cohérente étendant~$\Th{Q}$, dotée d'un prédicat de dérivabilité qui
satisfait ces conditions.

Il est instructif de suivre l'esquisse de
l'argument donnée par Gödel : le théorème en est
la conclusion frappante. Soit $!G_\Th{PA}$
l'!!{sentence} de Gödel construit dans la preuve du
\olref[1in]{thm:first-incompleteness}. Nous avons
montré : « Si $\Th{PA}$ est cohérente, elle ne
démontre pas $!G_\Th{PA}$. » En formalisant cette
implication \emph{dans} $\Th{PA}$, nous obtenons
une preuve de
\[
\OCon[\Th{PA}] \lif
\lnot\OProv[\Th{PA}](\gn{!G_\Th{PA}}).
\]
Supposons maintenant que $\Th{PA}$ démontre
$\OCon[\Th{PA}]$. Elle démontre alors
$\lnot\OProv[\Th{PA}](\gn{!G_\Th{PA}})$.
Mais, puisque $!G_\Th{PA}$ est un !!{sentence}
de Gödel, cette dernière formule équivaut à
$!G_\Th{PA}$. Ainsi, $\Th{PA}$ démontre
$!G_\Th{PA}$.

Or, si $\Th{PA}$ est cohérente, nous savons
qu'elle ne démontre pas $!G_\Th{PA}$!{} Elle ne
peut donc pas non plus démontrer
$\OCon[\Th{PA}]$.

Précisons l'argument. Prenons pour $!G_\Th{PA}$
l'!!{sentence} de Gödel de~$\Th{PA}$ et utilisons
les conditions de dérivabilité P1--P3 afin de
montrer que $\Th{PA}$ démontre
$\OCon[\Th{PA}]\lif !G_\Th{PA}$. Puisque
$\Th{PA}$ ne démontre pas $!G_\Th{PA}$ lorsqu'elle
est cohérente, elle ne démontre pas non plus
$\OCon[\Th{PA}]$. Voici les étapes de la preuve,
effectuées dans~$\Th{PA}$. Nous omettons les indices
$\Th{PA}$ pour alléger les notations et écrivons
$\OCon\ident\lnot\OProv(\gn{\lfalse})$.
\begin{align}
& !G \liff \lnot \OProv(\gn{!G}) \ollabel{G2-1}\\
& \qquad\text{$!G$ est une phrase de Gödel}\notag \\
& !G \lif \lnot \OProv(\gn{!G}) \ollabel{G2-2}\\
& \qquad\text{par \olref{G2-1}} \notag\\
& !G \lif
  (\OProv(\gn{!G}) \lif \lfalse) \ollabel{G2-3}\\
& \qquad\text{par \olref{G2-2} et la logique}\notag\\
& \OProv(\gn{
    !G \lif
    (\OProv(\gn{!G}) \lif \lfalse)
  }) \ollabel{G2-4}\\
& \qquad\text{par \olref{G2-3} et P1} \notag\\
& \OProv(\gn{!G}) \lif
  \OProv(\gn{
    (\OProv(\gn{!G}) \lif \lfalse)
    }) \ollabel{G2-5}\\
& \qquad\text{par \olref{G2-4} et P2} \notag\\
& \OProv(\gn{!G}) \lif (\OProv(\gn{\OProv(\gn{!G})}) \lif \OProv(\gn{\lfalse})) \ollabel{G2-6}\\
& \qquad\text{par \olref{G2-5}, P2 et la logique} \notag\\
& \OProv(\gn{!G}) \lif
  \OProv(\gn{\OProv(\gn{!G})}) \ollabel{G2-7}\\
& \qquad\text{par P3} \notag\\
& \OProv(\gn{!G}) \lif \OProv(\gn{\lfalse}) \ollabel{G2-8}\\
& \qquad\text{par \olref{G2-6}, \olref{G2-7} et la logique}\notag\\
& \OCon \lif \lnot \OProv(\gn{!G}) \ollabel{G2-9}\\
& \qquad\text{par contraposition de \olref{G2-8}}\notag \\
& \OCon \lif !G \notag\\
& \qquad\text{par \olref{G2-1}, \olref{G2-9} et la logique}\notag
\end{align}
Les recours à la logique ci-dessus n'emploient que
des faits élémentaires de logique propositionnelle.
Par exemple, \olref{G2-3} utilise
$\Proves\lnot !A\liff(!A\lif\lfalse)$, tandis
que \olref{G2-8} utilise
$!A\lif(!B\lif !C),\ !A\lif !B\Proves !A\lif !C$.
L'emploi de P2 en \olref{G2-5} et \olref{G2-6}
repose sur des instances de
$\OProv(\gn{!A\lif !B})\lif
(\OProv(\gn{!A})\lif\OProv(\gn{!B}))$.
Dans la première, $!A\ident !G$ et
$!B\ident\OProv(\gn{!G})\lif\lfalse$;
dans la seconde, $!A\ident\OProv(\gn{!G})$
et $!B\ident\lfalse$.

La version plus générale du second théorème
d'incomplétude est la suivante.

\begin{thm}
\ollabel{thm:second-incompleteness-gen}
Soit $\Th{T}$ une théorie cohérente, axiomatisée
et étendant~$\Th{Q}$, et soit
$\OProv[\Th{T}](y)$ une !!{formula} sans variable libre autre que~$y$, satisfaisant
les conditions de dérivabilité P1--P3 pour~$\Th{T}$.
Alors $\Th{T}$ ne démontre pas
$\OCon[\Th{T}]$, où
$\OCon[\Th{T}]\ident
\lnot\OProv[\Th{T}](\gn{\lfalse})$.
\end{thm}

\begin{prob}
Montrez que $\Th{PA}$ démontre
$!G_{\Th{PA}}\lif\OCon[\Th{PA}]$.
\end{prob}

\begin{digress}
La leçon à retenir est qu'aucune théorie cohérente
« raisonnable » des mathématiques ne peut démontrer
son propre énoncé de cohérence. Supposons qu'une
telle théorie mathématique cohérente~$\Th{T}$ contienne $\Th{Q}$
ainsi que les raisonnements « finitaires » de
Hilbert, quels qu'ils soient. La théorie entière
$\Th{T}$ ne peut démontrer son propre énoncé de
cohérence; à plus forte raison, sa partie finitaire
ne peut pas non plus le démontrer. En ce sens,
il ne peut exister de preuve finitaire de la
cohérence de « toutes les mathématiques ».

Le mot « finitaire » laisse une certaine marge
d'interprétation. Dans son article de 1931, Gödel
admettait que ce que nous pourrions juger
« finitaire » dépasse peut-être les mathématiques
que Hilbert souhaitait formaliser. Mais Gödel se
montrait charitable : de nos jours, il est difficile
de voir ce qui pourrait raisonnablement être dit
finitaire sans être formalisable, par exemple,
dans~$\Th{ZFC}$, la théorie des ensembles de
Zermelo--Fraenkel avec l'axiome du choix.
\end{digress}

\end{document}
